%!TEX root = ../main.tex

% ==================================================================
\section*{Scope and lecture outline}
\phantomsection
\addcontentsline{toc}{section}{Scope and lecture outline}

These lectures are an introduction to \textbf{toric methods for
Calabi--Yau compactifications and the string landscape} for
mathematicians.  The organising
question is concrete: how much of a four-dimensional effective theory,
and in particular of its vacuum energy, can be computed from explicit
algebraic and toric data?  The candidate de~Sitter construction of
\cite{McAllister:2024lnt} is the central worked example, but the first
goal is more basic: to explain why explicit Calabi--Yau geometry is
needed at all.

The basic toric vocabulary (fans, cones, divisors, the secondary
fan etc.) is used throughout without being re-derived in the live
lectures; it is collected in the appendices. By contrast, the
constructions that are central to the worked example, in particular
Batyrev's construction of Calabi--Yau hypersurfaces from reflexive
polytopes and the role of FRSTs, are re-derived at the level of a
stated theorem in Lecture~1 (see Theorem~\ref{thm:batyrev}). The
appendices collect the definitions and standard results from algebraic
geometry, convex geometry, toric geometry, mirror symmetry,
orientifolds, and computational workflows that are needed to read the
notes as a self-contained reference.

Sections marked by an asterisk are included for completeness and for
later reference, but they are not part of the intended blackboard
delivery. In the live lectures I will use them only as optional
background or as answers to questions.

\medskip
\noindent\textbf{Prerequisites.}
\begin{itemize}[nosep,leftmargin=*]
    \item Algebraic and complex geometry at the level of Huybrechts
    (\emph{Complex Geometry} \cite{Huybrechts:2005complex}) or
    Voisin (\emph{Hodge Theory and Complex Algebraic Geometry}
    \cite{Voisin:2002}).
    \item Toric geometry at the level of Cox--Little--Schenck
    \cite{cox-little-schenck:toric}: fans, divisors, intersection
    rings, the K\"ahler and Mori cones, the secondary fan, and
    reflexive polytopes. Batyrev's construction
    \cite{Batyrev:1993oya} and the Kreuzer--Skarke list
    \cite{Kreuzer:2000xy} are helpful background but are recalled in
    Lecture~1 (see \Cref{ssec:L1-toric}).
    \item Riemannian and Hermitian geometry: holonomy and the Hodge
    decomposition. Berger's classification is helpful background but
    is restated in Lecture~1 (see \Cref{thm:berger}).
    \item Basic physics vocabulary is helpful but not assumed: fields,
    actions, equations of motion, compactification, fluxes, branes,
    and effective field theory are introduced when first used.
\end{itemize}


\medskip
\noindent\textbf{Motivation.}
The motivation is two-fold. First, the notes explain how toric
geometry makes certain string compactifications explicit enough that
one can compute physical quantities rather than only discuss them
abstractly. The conceptual ingredients have long been established:
Calabi--Yau geometry, mirror symmetry, Picard--Fuchs equations,
Gopakumar--Vafa invariants, D-branes, and four-dimensional
supergravity have all been studied for decades. What has changed is
the feasible computational range. Many calculations that were once
limited to one- or few-modulus examples can now be attempted for the largest known Hodge numbers, i.e., compactifications with hundreds of moduli, provided one keeps careful control of
truncations and chamber structure.

Second, the notes are meant to identify bottlenecks. 
The concrete questions are geometric before they are
physical. Given a Calabi--Yau threefold $X$, how does one determine
its K\"ahler cone, effective cone etc.?
How does one compute minimal-volume representatives of divisor
classes, especially when holomorphic representatives exist only after
moving in the space of birational equivalence classes of $X$? Which effective divisors have
$h^\bullet(D,\mathcal{O}_D)=(1,0,0)$ and can therefore contribute to
the non-perturbative superpotential? How should one compare the Wall
data of hypersurfaces arising from different triangulations, and how
many inequivalent geometries are really present in the
Kreuzer--Skarke ensemble? 
The hope is that the lectures make these questions precise enough to
be useful to both the mathematics and string theory communities.

\medskip
\noindent\textbf{Lecture outline.}
The live blackboard narrative is deliberately narrower than the
written notes: compactification first selects Ricci-flat internal
geometries; the deformations of those geometries become
four-dimensional moduli; fluxes, sources and quantum corrections
generate a scalar potential $V$; and toric geometry supplies much
of the discrete data needed to write and analyse that potential. The
detailed explanation of how ten-dimensional fields reduce to
four-dimensional fields appears only after the Calabi--Yau moduli
have been introduced. In this sense the purpose is not to give a
general survey of string theory, but to explain why explicit
geometry is needed to compute its four-dimensional consequences.



There is intentional overlap with the TASI lectures
\cite{McAllister:2025qwq}. The TASI notes are the physics-first
reference for the Type IIB flux compactification, the leading-order
effective field theory, and the candidate de~Sitter construction.
The present notes reorganise that material for the ESI audience: the
blackboard route starts from compact Ricci-flat geometry, then
explains which toric computations supply the input to the
four-dimensional scalar potential. Whenever a statement is imported
from the TASI treatment rather than rederived here, the relevant TASI
section is cited explicitly.

\begin{enumerate}[leftmargin=*]
    \item \textbf{Lecture~1: The why and how of Calabi--Yau.}
    The first lecture introduces the minimal string theory input,
    explains why compactification leads to a
    Ricci-flat internal geometry, why Calabi--Yau threefolds are the
    computable and best understood class used in these notes, and how
    Batyrev's construction turns reflexive polytopes and FRSTs into
    explicit Calabi--Yau hypersurfaces in toric varieties. The emphasis
    is pragmatic: what toric methods compute reliably today
    (Hodge data, intersection numbers, cones, divisor topology), and
    what remains beyond current algorithms.
    Along the way we preview where these data enter physics: flux
    superpotentials, instanton corrections, divisor zero-mode tests,
    and K\"ahler-cone navigation.

    \item \textbf{Lecture~2: Periods and GV invariants from toric data.}
    The second lecture introduces the minimal Type IIB/O3--O7
    dictionary needed to write the flux superpotential and then
    focuses on the main computational problem:
    \begin{equation}
        W_{\rm flux}=\int_X\Omega\wedge G_3 .
    \end{equation}
    Fluxes supply the integral coefficients; periods supply the
    functions. The lecture explains how mirror symmetry, toric
    hypersurfaces, the mirror map, and genus-zero GV invariants make
    these periods computable. It then uses PFVs as the final
    application: special integral fluxes make the polynomial LCS
    contribution flat, while the GV-controlled instanton terms lift
    the remaining direction and produce a computable small value of
    $W_0$. Non-perturbative K\"ahler-modulus stabilisation and the
    candidate-dS pipeline are deferred to Lecture~3.

    \item \textbf{Lecture~3: Vacuum energies from toric geometry.}
    The third lecture assembles the leading-order effective
    theory for Type IIB O3/O7 orientifold compactifications with flux,
    non-perturbative K\"ahler-modulus stabilisation, and a
    redshifted antibrane uplift. It introduces
    geometric phases, flop transitions, vex chambers, and the
    extended K\"ahler cone,
    then explains how chamber-aware algorithms, including
    \texttt{fanroots} and \texttt{JAXVacua}, locate
    K\"ahler moduli critical points. The remaining ingredients are the
    Klebanov--Strassler throat, KPV anti-D3 uplift, the explicit
    toric AdS precursor and candidate-dS recipes, and the control
    checks: tadpoles, GV convergence, higher-order corrections,
    backreaction, and quantisation issues. The list is retained
    because these checks distinguish the leading-order candidate
    construction from a fully controlled vacuum theorem.
\end{enumerate}

% ==================================================================
\setcounter{section}{0}
